% ~3 min
\begin{frame}
  \frametitle{Why Hylomorphisms in Rocq?}
  \begin{itemize}
    \item Haskell: the schemes, but no termination.
    \item Rocq: termination, but these definitions are rejected.
    \item We want: a \embf{library}, \embf{extracted code}, and the
      \embf{laws as tactics}.
  \end{itemize}
  \speakernote{Library: termination is proved once per coalgebra, not per
    function. Extracted code: OCaml or Haskell with a termination
    guarantee. Laws as tactics: program calculation and optimisation by
    rewriting.}
\end{frame}

\begin{frame}
  \frametitle{Challenges}
  \begin{enumerate}
    \item \textbf{No axioms}
    \item \textbf{Clean extracted code}
    \item \textbf{Non-structural recursion}
  \end{enumerate}
  \uncover<2->{%
    Solutions:
    \begin{enumerate}
      \item<2-> Setoids; decidable validity.
      \item<3-> One type of positions, not a family.
      \item<4-> Recursion on a termination proof.
    \end{enumerate}}
  \uncover<5->{\textbf{Since the paper:} usability.}
  \speakernote{The axioms to avoid: functional extensionality, axiom K,
    heterogeneous equality. Clean code: what a programmer would write,
    no Obj.magic. Divide-and-conquer recursion is not structural. The
    solutions are Parts I and II. Since the paper the hard part has been
    usability (instance search, rewriting, inlining): end of Part III.}
\end{frame}

\begin{frame}
  \frametitle{Roadmap}
  \begin{center}
    \Large
    \begin{sticky}
      \begin{tabular}{@{}rl}
        \textbf{\color{gray}Part I:} & Extractable Containers \\
        \textbf{\color{gray}Part II:} & Recursive Coalgebras \\
        \textbf{\color{gray}Part III:} & Calculation and Extraction \\
        \textbf{\color{gray}Part IV:} & Dynamic Programming \\
        \textbf{\color{gray}Part V:} & Monoid Algebras \\
      \end{tabular}
    \end{sticky}
  \end{center}
  \speakernote{Parts I and II are the ITP paper; Parts III to V include
    newer material.}
\end{frame}
